\documentclass[10pt,a4paper]{amsart}
\usepackage[margin=19mm]{geometry}
\usepackage{amsmath,amssymb,mathtools}
\usepackage[scr=rsfs,cal=euler]{mathalfa}
\usepackage{xfrac}
\usepackage[unicode,hidelinks,bookmarks=false]{hyperref}
\newcommand{\ie}{i.e.\ }
\newcommand{\cf}{cf.\ }
\newcommand{\resp}{respectively\ }
\newcommand{\Subsection}{\subsection}
\providecommand{\nobreakdash}{\nobreak\hbox{-}}
\setlength{\emergencystretch}{2em}
\multlinegap0pt
\usepackage{bm}
\usepackage{bbm}
\usepackage{dsfont}
\usepackage{xspace}
\usepackage{cancel}
\usepackage{mathtools}
\usepackage{amsthm}
\makeatletter
\@ifundefined{thm}{\newtheorem{thm}{Thm}}{}
\@ifundefined{lem}{\newtheorem{lem}[thm]{Lemma}}{}
\@ifundefined{prop}{\newtheorem{prop}[thm]{Proposition}}{}
\makeatother
\let\phi=\varphi
\newcommand{\D}{\mathbb{D}}
\newcommand{\T}{\bm{T}}
\newcommand{\cla}{\mathcal{A}}
\newcommand{\clb}{\mathcal{B}}
\newcommand{\cle}{\mathcal{E}}
\newcommand{\cln}{\mathcal{N}}
\newcommand{\clr}{\mathcal{R}}
\title[Characteristic function of a power partial isometry]{Characteristic function of a power partial isometry}
\author{Kritika Babbar, Amit Maji}
\date{Source: arXiv:2505.07824v3 (CC BY 4.0)}
\begin{document}
\maketitle

\noindent\emph{Source and licence.} Characteristic function of a power partial isometry. Version arXiv:2505.07824v3 (CC BY 4.0), distributed under the Creative Commons Attribution 4.0 International licence, as recorded in the arXiv licence metadata for this exact version and shown on the abstract page https://arxiv.org/abs/2505.07824v3. Full rights notice: licenses/arxiv-2505.07824-CC-BY-4.0.txt.\par\smallskip

\noindent\textit{Editorial scope.} Counted unit: the Prop whose source label is \texttt{NC} in the article (source0.tex lines 1086--1094, source label \texttt{NC}), delivered together with the article's own complete original proof of it (source0.tex lines 1096--1140, one \texttt{proof} environment of 45 lines). No mathematics is written, repaired or re-derived: statement and proof are the article's own text line for line. Required context retained uncounted: NC1 (lines 1059--1061). Every definition, display and label of those lines is retained unchanged. Omitted from this package and not redistributed: 1--1058, 1062--1085, 1095--1095, 1141--1646. Nothing inside a retained excerpt is omitted or altered: every retained body is delivered exactly as the article's lines are written, so no omission receipt of any kind accompanies this package. Citations: the retained excerpts carry no citation command at all, so no bibliography entry of the article is transcribed and no reference is redistributed. Reference adaptation: none was needed and none was made; every reference inside the delivered unit resolves inside it, so no unresolved reference or citation is produced. Printed slips kept literal: no printed word, symbol, hypothesis or number of the counted statement or of its proof was altered. Layout and class: the article's own class is \texttt{amsart} with packages including amssymb, amsmath, amsfonts, latexsym, bm, xy, amscd, setspace, mathrsfs, hyperref; the wrapper is the lane's pinned \texttt{amsart} editorial wrapper at 10pt on A4, which restates the theorem-like environments the retained text needs and adds the article's own macro definitions that the retained text needs (\texttt{\textbackslash D}, \texttt{\textbackslash T}, \texttt{\textbackslash cla}, \texttt{\textbackslash clb}, \texttt{\textbackslash cle}, \texttt{\textbackslash cln}, \texttt{\textbackslash clr}, \texttt{\textbackslash phi}), with extra wrapper packages (bm, bbm, dsfont, xspace, cancel, mathtools). The delivered numbering is the unit's own. Assets: no retained span contains a figure environment, a graphics-inclusion command, a PDF-page inclusion, a table or a TikZ picture; no image file is copied into this package, the package declares no asset dependency and no image file is redistributed with it. Licence: version arXiv:2505.07824v3 of this article is distributed under the Creative Commons Attribution 4.0 International licence, as recorded in the arXiv licence metadata for this exact version and shown on the abstract page https://arxiv.org/abs/2505.07824v3. A full rights notice is delivered at licenses/arxiv-2505.07824-CC-BY-4.0.txt. Characters: every retained excerpt is pure ASCII, so the delivered engine, XeLaTeX with its default Latin Modern text font, reports no missing character.
\begin{lem}\label{NC1}
Let $\Phi \in L^\infty_{\clb(\cle,\cle_*)}$ be such that $T_\Phi: H^2_\cle(\D) \rightarrow H^2_{\cle_*}(\D)$ is a nonzero partially isometric Toeplitz operator. Then $\Phi(e^{it})$ is a partial isometry a.e.\ on $\T$ and $\Phi \in \cla.$
\end{lem}

\begin{prop}\label{NC}
Let $\Phi \in \cla$ be such that $T_\Phi: H^2_\cle(\D) \rightarrow H^2_{\cle_*}(\D)$ is a nonzero partially isometric Toeplitz operator. Then $\Phi$ satisfies the following conditions:
\begin{enumerate}
\item  $\Phi_+(e^{it})^*\Phi_+(e^{it}) $ and $\Phi_-(e^{it})\Phi_-(e^{it})^*$
 are  operator-valued constant functions a.e.\ on $\T$, where $\Phi_+ \in H^\infty_{\clb(\cle,\cle_*)}$ and $\Phi_-^* \in H^\infty_{\clb(\cle_*,\cle)}$ are contractive analytic functions. \label{con1} 
\item $\hat{\Phi}(m)^*\hat{\Phi}(-n)=0_\cle$ and $\hat{\Phi}(-n)\hat{\Phi}(m)^*=0_{\cle_*}$ for all $m,n \geq 1$.
\label{con2}
\end{enumerate}
\end{prop}

\begin{proof}
Suppose $\Phi \in \cla$ such that $T_\Phi: H^2_\cle(\D) \rightarrow H^2_{\cle_*}(\D)$ is a nonzero partial isometric Toeplitz operator.
Write
\begin{equation}\label{eq 6}
\Phi= \sum\limits_{m=1}^\infty  \phi_{-m}\bar{z}^m +\sum\limits_{m=0}^\infty \phi_mz^{m}.
\end{equation}


From the proof of Lemma \ref{NC1}, $\Phi_+=\Theta\Psi(0)^*$ and $\Phi_-=\Theta(0)\Psi^*$ for operator valued inner functions $\Theta$ and $\Psi$. Observe that
\[
(\Theta(e^{it}) \psi_0^*)^*(\Theta(e^{it}) \psi_0^*)=\psi_0\psi_0^* \,\, \quad(\text{a.e.\ on } \T).
\]
Thus we get $\Phi_+(e^{it})^*\Phi_+(e^{it}) $ is a constant positive operator a.e.\ on $\T$. Similarly,  $\Phi_-(e^{it})\Phi_-(e^{it})^*=(\theta_0 \Psi(e^{it})^*)(\theta_0 \Psi(e^{it})^*)^* = \theta_0 \theta_0^*$ is a constant positive operator a.e.\ on 
$\T$. 

Now let $\phi_{-m_0}$ be nonzero for some $m_0 \geq 1$. For $\eta \in \cle$, define a function $f\in H^2_{\cle_*}(\D)$ by 
\[
f(z)=(T_\Phi M_z-M_zT_\Phi)\left(\eta z^{m_0-1}\right).
\]
Then 
\begin{align*}
f(e^{it})&=T_\Phi(\eta e^{im_0t})-e^{it}T_\Phi(\eta e^{i(m_0-1)t})\\
&=P_+\left(\sum\limits_{m=1}^\infty \phi_{-m}(\eta)e^{i(-m+m_0)t}\right)+\sum\limits_{m=0}^\infty \phi_m(\eta)e^{i(m+m_0)t} \\
& \quad -e^{it}P_+\left(\sum\limits_{m=1}^\infty \phi_{-m}(\eta)e^{i(-m+m_0-1)t}\right) - e^{it}\sum\limits_{m=0}^\infty \phi_{m}(\eta)e^{i(m+m_0-1)t}\\
&=\sum\limits_{m=1}^{m_0} \phi_{-m}(\eta)e^{i(-m+m_0)t}-e^{it}\left(\sum\limits_{m=1}^{m_0-1}\phi_{-m}(\eta)e^{i(-m+m_0-1)t}\right)\\
&=\phi_{-m_0}(\eta).
\end{align*}
Since $T_\Phi$ is a partial isometry, $T_\Phi^*$  is also a partial isometry. 
Hence $[\cln(T_\Phi^*)]^\perp=\clr(T_\Phi)$ is $M_z$-invariant which yields $f \in \clr(T_\Phi)$.
Also note that 
\[
\|f\|=\|T_\Phi^*f\|=\|P_+L_\Phi ^*f\|\leq \|L_\Phi^* f \|\leq \|f\|.
\]
Therefore, $\|P_+L_\Phi ^*f \|=\|L_\Phi^* f \|$ which implies $P_+ L_\Phi^* (f)=L_\Phi^*f$. Hence $L_\Phi^* f \in H^2_\cle(\D)$. Then
\[
\sum\limits_{m=1}^\infty \phi_{-m}^*\phi_{-m_0}(\eta)e^{imt}+\sum\limits_{m=0}^\infty \phi_m^*\phi_{-m_0}(\eta)e^{-imt} = \Phi(e^{it})^*f(e^{it})\in H^2_\cle(\D)
\]
if and only if 
\[
\phi_m^*\phi_{-m_0}\eta=0 \quad \mbox{for all} ~~ m \geq 1.
\]
Since $m_0 \geq 1$ is arbitrary such that $\phi_{-m_0} \neq 0$, we have $\phi_m^*\phi_{-n}=0 $ for all $m,n \geq 1$. Note that if $\phi_{-m}=0$ for all $m \geq 1$, then this condition is trivially true. 

Again, since $T_\Phi^*$ is also a partial isometry, on replacing $\Phi$ by $\Phi^*$, we obtain $\phi_{-n}\phi_m^*=0$ for all $m ,n \geq 1$. This finishes the proof.
\end{proof}

\end{document}
